\documentclass[a4paper]{article}
% Español (México) con XeLaTeX: fontspec para fuentes Unicode, polyglossia
% para la división silábica y los nombres del documento en la variante mexicana.
\usepackage{fontspec}
\usepackage{polyglossia}
\setdefaultlanguage[variant=mexican]{spanish}
% Los nombres chinos van como «pinyin (original)»: xeCJK da a los caracteres
% Han su propia fuente; sin ella XeLaTeX los omite con un aviso.
% FandolSong no cubre algunos caracteres de nombres (U+73FA); WenQuanYi Zen Hei sí.
\usepackage[AutoFallBack=true]{xeCJK}
\setCJKmainfont{FandolSong-Regular.otf}
\setCJKfallbackfamilyfont{\CJKrmdefault}{WenQuanYi Zen Hei}
% polyglossia no traduce los nombres de listings.
\AtBeginDocument{\providecommand{\lstlistingname}{}\renewcommand{\lstlistingname}{Listado}%
  \providecommand{\lstlistlistingname}{}\renewcommand{\lstlistlistingname}{Índice de listados}}
\usepackage{amsmath, amssymb}
\usepackage{graphicx}
\usepackage[margin=2.35cm]{geometry}
\usepackage[most]{tcolorbox}
\usepackage{booktabs}
\usepackage{tabularx}
\usepackage{array}
\usepackage{float}
\usepackage{xcolor}
\usepackage{hyperref}
\usepackage{fancyhdr}
\setCJKmainfont[AutoFakeBold=2.4]{AR PL UMing CN}
\setCJKsansfont[AutoFakeBold=2.4]{Droid Sans Fallback}
\setCJKmonofont[AutoFakeBold=2.4]{Droid Sans Fallback}
\hypersetup{colorlinks=true, linkcolor=blue!55!black, urlcolor=blue!60!black}
\graphicspath{{./}{figures/}}
\setlength{\parskip}{0.55em}
\setlength{\parindent}{2em}
\setlength{\emergencystretch}{3em}
\renewcommand{\arraystretch}{1.18}
\newtcolorbox{knowledgebox}[1]{enhanced, colback=blue!5!white, colframe=blue!70!black, colbacktitle=blue!70!black, coltitle=white, fonttitle=\bfseries, title={#1}, sharp corners, breakable}
\newtcolorbox{importantbox}[1]{enhanced, colback=yellow!10!white, colframe=yellow!75!black, colbacktitle=yellow!75!black, coltitle=black, fonttitle=\bfseries, title={#1}, sharp corners, breakable}
\newtcolorbox{warningbox}[1]{enhanced, colback=red!5!white, colframe=red!70!black, colbacktitle=red!70!black, coltitle=white, fonttitle=\bfseries, title={#1}, sharp corners, breakable}
\pagestyle{fancy}
\fancyhf{}
\fancyfoot[C]{\thepage}
\fancyfoot[R]{\footnotesize\color{gray}\href{https://github.com/hqhq1025/ai-course-notes}{AI Course Notes}}
\renewcommand{\headrulewidth}{0pt}
\renewcommand{\footrulewidth}{0pt}
\newcommand{\videourl}{https://www.youtube.com/watch?v=Puptr04av5g}
\newcommand{\repourl}{https://github.com/hqhq1025/ai-course-notes}

\begin{document}
\begin{titlepage}
\centering
{\Large Notas didácticas de una entrevista a fondo\par}
\vspace{1.0cm}
{\huge\bfseries Inteligencia corporizada: su historia en los márgenes de la academia, el problema de los datos y la ruta hacia la productividad}\par
\vspace{0.5cm}
{\Large Zhang Xiaojun conversa con Wang He (王鹤), fundador y CTO de Yinhe Tongyong (银河通用)\par}
\vspace{0.35cm}
{\large Elaboradas a partir del video público, la descripción del video, la transcripción, la estructura de capítulos y diagramas conceptuales propios del Zhang Xiaojun Podcast\par}
\vspace{0.3cm}
{\large 2025-07-02\par}
\vspace{0.85cm}
\includegraphics[width=0.88\textwidth,height=0.42\textheight,keepaspectratio]{cover.jpg}\par
\vfill
\begin{tcolorbox}[width=0.92\textwidth, colback=black!2!white, colframe=black!55, sharp corners]
\textbf{Título del video}: Conversación con Wang He: la historia de la inteligencia corporizada en los márgenes de la academia y el desorden provocado por el bombardeo de capital (con subtítulos)\par
\textbf{Canal del video}: Zhang Xiaojun Podcast\par
\textbf{Duración del video}: 02:38:11\par
\textbf{Enlace del video}: \href{\videourl}{\nolinkurl{\videourl}}\par
\textbf{Nota sobre los materiales}: YouTube no ofrece subtítulos descargables; el texto se elaboró a partir de la descripción del video, la información de capítulos, la transcripción de faster-whisper y figuras didácticas propias. Las tomas fijas de la entrevista no se repiten en el cuerpo del texto.\par
\textbf{Página del proyecto}: \href{\repourl}{\nolinkurl{\repourl}}\par
\end{tcolorbox}
\end{titlepage}

\tableofcontents
\newpage
\section{Introducción: del furor por el concepto de vuelta a la productividad}

Esta sección establece primero la manera de leer el episodio completo. Wang He (王鹤) es a la vez profesor asistente de la Universidad de Pekín y fundador y CTO de Yinhe Tongyong (银河通用); su relato abarca al mismo tiempo la historia académica, las rutas tecnológicas, las decisiones de emprendimiento y la ética de la industria. Más que discutir solo «si los robots llegarán pronto a los hogares», las preguntas más importantes de este episodio son: por qué la inteligencia corporizada pasó de ser una línea marginal de la visión por computadora a convertirse en la favorita del capital, por qué el lenguaje no es la única puerta de entrada a la inteligencia, por qué los datos para VLM/VLA son más escasos que los datos para LLM y por qué, en última instancia, la industria debe demostrar su valor con productividad y no con demos.

Este episodio complementa la entrevista con Xie Chen (谢晨) del EP109. El EP109 explica «cómo se fabrican los datos» desde la infraestructura de simulación y de datos sintéticos; el EP106 explica, desde las escuelas académicas, el ciclo cerrado visión-acción y la lógica de implantación de una empresa, «por qué se necesita este tipo de datos y qué narrativas comerciales dañan a la industria». En conjunto, los dos episodios permiten articular la inteligencia corporizada como una cadena completa que va del concepto en los artículos académicos a las rutas de modelos, la recipe de datos, las restricciones de hardware y el ROI en escenarios reales.

\begin{importantbox}{Tesis central del episodio}
La inteligencia corporizada no es algo tan simple como «meter un modelo grande en un robot». En primer lugar, exige que la percepción, la acción y la retroalimentación del entorno formen un ciclo cerrado; en segundo lugar, exige que el hardware, los datos, los modelos y los escenarios avancen en espiral ascendente; por último, debe demostrar su valor con productividad escalable, y no con montos de financiamiento, demos de teleoperación o narrativas vagas sobre robots de propósito general.
\end{importantbox}

\begin{knowledgebox}{Criterio visual}
Este video es una entrevista con encuadre fijo, sin slides, pizarrón ni demostraciones de producto. El texto usa la portada solo para identificar la fuente; todas las imágenes del texto son diagramas conceptuales de elaboración propia, que sirven para explicar el lenguaje y la inteligencia, la genealogía académica de la inteligencia corporizada, el Perception-Action Loop, la recipe de datos, la productividad como producto y los excesos del capital.
\end{knowledgebox}

\subsection{Resumen de la sección}

El hilo conductor del EP106 es que la inteligencia corporizada pasó del margen académico de la visión al centro de la industria, pero el verdadero criterio no es qué tanto entusiasmo despierta la historia, sino si se logra poner en marcha al mismo tiempo el ciclo cerrado percepción-acción, el volante de datos y la productividad comercial.
\section{El lenguaje no es la esencia de la inteligencia, sino un salto cualitativo}

Esta sección parte del concepto que más se presta a confusión. En la entrevista, Wang He (王鹤) define la inteligencia de un modo cercano a «la capacidad de responder al entorno según la situación». Esta definición no pone el lenguaje en el centro de forma deliberada, porque muchos seres vivos carecen de un lenguaje de orden superior y aun así muestran inteligencia en su entorno; la visión, el tacto, el sonido y la propiocepción pueden ser los sensor con los que un agente inteligente percibe el mundo. La importancia del lenguaje reside en que permite a los seres humanos comprimir la experiencia, difundir el conocimiento, formar conceptos abstractos y aprender de manera colaborativa; por eso es un salto cualitativo, y no la única fuente de la inteligencia.

Este juicio incide directamente en la orientación de la inteligencia corporizada. Si se equipara la inteligencia con el lenguaje, se tiende a pensar que los LLM ya resolvieron la mayor parte de los problemas; si la inteligencia se entiende como la capacidad de responder al entorno, se advierte que el robot debe volver a la percepción, la acción y la retroalimentación. Un modelo de lenguaje puede aportar sentido común y descomponer tareas, pero el robot tiene que actuar en el mundo físico, donde un fallo altera el entorno y tiene costos; esto exige otro tipo de datos y otra forma de evaluación.

\begin{figure}[H]
\centering
\includegraphics[width=0.94\textwidth]{language-intelligence-jump.png}
\caption{El lenguaje no es la esencia de la inteligencia: es un salto cualitativo de la inteligencia humana, y no la única vía de acceso a la inteligencia. Diagrama conceptual propio, elaborado a partir de la conversación de 00:05:58--00:25:08.}
\end{figure}

\begin{knowledgebox}{Lectura de la figura: el lenguaje potencia la inteligencia, pero no sustituye al cuerpo}
A la izquierda, intelligence pone el acento en el entorno, la respuesta y la retroalimentación; a la derecha, language lo pone en la compresión, la difusión y la colaboración. Al leer esta figura, no hay que ver ambos lados como opuestos, sino como una relación jerárquica: el cuerpo y la percepción proporcionan el contacto con el entorno, y el lenguaje permite que la experiencia se transmita entre individuos. El problema de la inteligencia corporizada es que no basta con que el lado del lenguaje sea fuerte; el lado físico también debe cerrar el ciclo.
\end{knowledgebox}

\subsection{La visión es un sensor potente, no una inteligencia completa}

Esta subsección distingue con más detalle entre visión e inteligencia. La visión es un sensor muy potente, que apareció antes que muchas otras modalidades perceptivas en la evolución de la inteligencia de los animales superiores; sin embargo, el reconocimiento puramente visual suele ser passive perception, es decir, dada una imagen, determinar su categoría, segmentar regiones o estimar posiciones. Una vez concluido el reconocimiento, el entorno no cambia porque uno diga «esto es un gato», ni proporciona retroalimentación para el siguiente paso. Esa es la diferencia entre las tareas de visión basadas en internet y las tareas de inteligencia corporizada.

Wang He señala que el campo de la visión por computadora acumuló durante años una enorme capacidad gracias a tareas como ImageNet, el reconocimiento facial y la segmentación semántica, pero estas tareas estaban impulsadas en su mayoría por datos de internet y anotación manual. Cuando la competencia en esas tareas llegó a cierto grado de saturación, los investigadores empezaron a buscar el siguiente paso: hacer que la visión fuera activa, que el agent emprendiera acciones, que las acciones modificaran las observaciones y que las nuevas observaciones guiaran el siguiente paso. De ahí surge el Perception-Action Loop.

\begin{warningbox}{No hay que equiparar la capacidad de reconocimiento con la capacidad corporizada}
Poder reconocer una taza, una silla y una habitación no equivale a poder desplazarse junto a la silla, tomar la taza, abrir la puerta de un gabinete u ordenar un anaquel. La inteligencia corporizada requiere acción y retroalimentación del entorno, no solo asignar etiquetas a imágenes estáticas.
\end{warningbox}

\subsection{Resumen de la sección}

El lenguaje produjo un salto en la inteligencia humana, pero la inteligencia corporizada no puede depender solo del lenguaje. El robot necesita enlazar la visión, la acción y la retroalimentación del entorno en un ciclo cerrado; este es también el motor principal del giro de la investigación en visión por computadora hacia la inteligencia corporizada.
\section{Historia académica marginal de la inteligencia corporizada}

La sección anterior explicó los límites entre el lenguaje y la visión; esta sección entra en la historia académica. Wang He (王鹤) subraya que «inteligencia corporizada» y «robótica» no fueron, en su origen, un mismo relato sostenido por un mismo grupo de investigadores. La inteligencia corporizada surgió sobre todo de la búsqueda, por parte de los investigadores de visión por computadora, del problema de la siguiente etapa: pasar de la passive perception a agentes capaces de actuar, de cambiar lo que observan y de aprender en lazo cerrado. Los investigadores de la tradición robótica, en cambio, se han ocupado durante mucho tiempo del control, la estructura mecánica, la optimización, los escenarios industriales y el hardware concreto.

Esta diferencia explica por qué la tarea más importante en los inicios de la inteligencia corporizada no fue la manipulación diestra, sino la navegación. La navegación permite a los investigadores de visión permanecer en una zona de confort relativamente conocida: el agent cambia su propia posición, la cámara capta una imagen nueva y el entorno no tiene que modificarse. Ya cuenta con un ciclo de percepción-acción, pero es más fácil de estandarizar que el agarre, la apertura de puertas, el contacto y el control de fuerza. Por ello, tareas y plataformas como Object Goal Navigation, Point Goal Navigation y Habitat se convirtieron en el centro del relato inicial.

\begin{figure}[H]
\centering
\includegraphics[width=0.96\textwidth]{embodied-ai-origin.png}
\caption{Historia académica marginal de la inteligencia corporizada: de la visión, la navegación y el ciclo de percepción y acción hacia la corriente principal. Diagrama conceptual propio, elaborado a partir de la conversación entre 00:25:08 y 00:41:15.}
\end{figure}

\begin{knowledgebox}{Lectura de la figura: la corriente académica no surgió de manera natural de la «industria robótica»}
En la figura, el recorrido de Vision a Navigation y después a Perception y Action muestra que la inteligencia corporizada fue, en su origen, una extensión de la «siguiente generación de tareas de visión» propuesta por investigadores de visión. Más adelante incorporó robot learning, el agarre, el aprendizaje por refuerzo y los modelos grandes, pero su relato central fue al principio este: la visión no puede limitarse a reconocer de forma pasiva, sino que debe integrarse en un lazo cerrado de acción.
\end{knowledgebox}

\subsection{Perception-Action Loop}

Esta subsección expone con claridad el paradigma de fondo de la inteligencia corporizada. El Perception-Action Loop consiste en «decidir primero la acción a partir de la percepción; la acción modifica el entorno o la propia posición, y la nueva percepción guía el paso siguiente». La navegación es el ejemplo más sencillo: el robot avanza, cambian sus coordenadas y el ángulo de la cámara, y en consecuencia ve un entorno nuevo. Las tareas de manipulación son más complejas: después de tomar una botella, cambian la posición de la botella, el estado de lo que se sostiene en la mano y las acciones posibles a continuación.

En comparación con el agarre de un solo paso, el lazo cerrado refleja mejor el valor de investigación de la inteligencia corporizada. Muchos de los primeros métodos de agarre convertían el problema en razonamiento sobre nubes de puntos o geometría: observar el objeto una vez, predecir el punto de agarre, ejecutar una sola vez y dar la tarea por terminada. Aunque esto ocurre en un robot, sigue pareciéndose a una tarea de reconocimiento de patrones tridimensionales. Lo que persigue la inteligencia corporizada es el lazo cerrado de varios pasos: cada acción modifica el estado del mundo, y el modelo debe volver a observar, volver a planificar y seguir ejecutando.

\begin{figure}[H]
\centering
\includegraphics[width=0.96\textwidth]{perception-action-loop.png}
\caption{Perception-Action Loop: el relato central de la inteligencia corporizada es el lazo cerrado de percepción, acción y retroalimentación. Diagrama conceptual propio, elaborado a partir de la conversación entre 00:25:08 y 00:41:15.}
\end{figure}

\begin{knowledgebox}{Lectura de la figura: el lazo cerrado convierte la visión de clasificación en acción}
De Observe a Infer, y después a Act, Feedback y Adapt, el cambio decisivo es que el entorno ya no es solo una imagen de entrada, sino que genera estados nuevos como consecuencia de las acciones. Al leer la figura conviene prestar atención al último paso, Adapt: la inteligencia corporizada no es una predicción única, sino la corrección continua de la política mediante la retroalimentación.
\end{knowledgebox}

\begin{importantbox}{Hitos representativos}
Wang He menciona que Li Feifei (李飞飞) calificó la inteligencia corporizada como una de las tres estrellas polares del futuro de la visión por computadora, y que Jensen Huang (黄仁勋), en eventos de NVIDIA, ha insistido en que la próxima generación de IA es la embodied AI. Este tipo de declaraciones llevó una dirección académica marginal hacia la corriente principal e hizo que las comunidades de visión, robótica, aprendizaje por refuerzo y modelos grandes comenzaran a converger en torno a un mismo término.
\end{importantbox}

\subsection{Resumen de la sección}

La historia académica de la inteligencia corporizada es una migración hacia una «visión activa»: del reconocimiento de imágenes a la navegación, del razonamiento de un solo paso a la acción en lazo cerrado, y después la incorporación gradual de la manipulation, el robot learning y la capacidad de los modelos grandes para tareas generales.
\section{La trayectoria académica de Wang He (王鹤): del video a los datos sintéticos}

La sección anterior trató sobre la comunidad académica; esta trata sobre la trayectoria personal de Wang He, porque sus proyectos de doctorado anticiparon casi por completo las dos líneas principales que después seguiría la inteligencia corpórea: aprender relaciones de interacción a partir de videos de personas y resolver con datos sintéticos la escasez de anotaciones reales. El primer trabajo consistió en aprender, a partir de videos de personas, interacciones de varios pasos entre personas y objetos, para generar animaciones o accionar una demo sencilla de una taza inteligente. En esencia, debía aprender las relaciones causales entre el estado, la acción, el cambio de estado del objeto y la acción probable siguiente.

Esto ocurrió muy pronto, en 2016. En ese momento no existían los grandes modelos multimodales de hoy ni datos VLA consolidados. Para saber si la taza estaba vacía o llena, si la tapa de la botella estaba abierta o cerrada, si la mano estaba vacía o sostenía un objeto, y dónde estaban los límites de cada segmento de acción, había que dividir el problema en varios modelos, con anotación manual y un sistema complejo. Wang He concluyó más tarde que, aunque aprender un modelo del mundo únicamente a partir de video todavía no es hoy la técnica que más impulsa la inteligencia corpórea, ese proyecto le hizo ver muy pronto que el problema de fondo no es el reconocimiento, sino aprender, a partir de la interacción, cómo cambia el estado del mundo.

\begin{figure}[H]
\centering
\includegraphics[width=0.96\textwidth]{academic-path.png}
\caption{Trayectoria académica de Wang He: de los videos de personas y los modelos del mundo al objetivo de investigación en robots domésticos. Diagrama conceptual propio, elaborado a partir de la conversación de 00:41:15--01:25:08.}
\end{figure}

\begin{knowledgebox}{Lectura de la figura: dos proyectos de doctorado que corresponden a dos líneas principales de la inteligencia corpórea}
La primera línea consiste en comprender, a partir de video, las interacciones causales de varios pasos entre personas y objetos, algo cercano al problema actual del world model; la segunda es la estimación de pose de objetos a nivel de categoría, que usa datos sintéticos o de realidad mixta para suplir la falta de anotaciones reales. Al leer la figura, conviene verlas como dos caras del mismo objetivo: comprender el estado de la interacción y obtener a bajo costo datos con los que se pueda entrenar.
\end{knowledgebox}

\subsection{Del modelado físico a la investigación en IA}

Esta subsección explica por qué Wang He pudo abstraer interacciones complejas en modelos. Su formación de licenciatura y de los primeros años de doctorado se inclinaba hacia la física, los dispositivos y el modelado matemático; estaba acostumbrado a ajustar modelos a partir de datos experimentales y luego usarlos para predecir casos nuevos. Más adelante se dio cuenta de que, en el clean room de semiconductores, la iteración experimental era lenta y el control deficiente, lo cual no correspondía con su ritmo de pensamiento, así que se orientó hacia la IA, la graficación por computadora y la interacción física. Esta experiencia muestra que la inteligencia corpórea no consiste solo en escribir código: también exige comprender los estados de los objetos, la causalidad de las acciones y los cambios físicos.

Lo más difícil del primer proyecto de doctorado no fue una arquitectura de red en particular, sino formulate research problem: qué estados había que extraer, qué acciones modificaban qué estados de los objetos y cómo encadenar esos cambios en un modelo causal temporal. Wang He mencionó que, cuando todavía no conocía la expresión, después popular, Perception-Action Loop, ya le había dibujado a su asesor un diagrama de state-action-change world-state. Ese es precisamente el ciclo de base que más tarde adoptaría la inteligencia corpórea.

\begin{knowledgebox}{Términos clave: modelado de la interacción}
\begin{tabularx}{\textwidth}{p{0.20\textwidth}p{0.34\textwidth}X}
\toprule
Término & Significado & Función en este episodio \\
\midrule
Object State & Estado del objeto, como vacío/lleno, abierto/cerrado, posición/orientación & Determina si la acción siguiente es viable. \\
Action Grammar & Método tradicional que describe con reglas el orden de las acciones & Es interpretable, pero tiene una capacidad limitada de generalización y de aprovechar los datos. \\
World Model & Modelo que predice cómo las acciones cambian el estado del mundo & Uno de los objetivos de largo plazo de la inteligencia corpórea. \\
Perception-Action Loop & Ciclo de percepción, acción y retroalimentación & Convierte el reconocimiento estático en inteligencia de ciclo cerrado. \\
\bottomrule
\end{tabularx}
\end{knowledgebox}

\subsection{Estimación de pose a nivel de categoría y datos sintéticos}

El segundo trabajo está más cerca de la vía actual de datos para la inteligencia corpórea. La 6D pose estimation tradicional suele dirigirse a una instancia conocida: primero se define para el objeto un origen de coordenadas y una orientación estándar, y luego se predicen su rotación y su traslación respecto de ese sistema de coordenadas. El problema que planteó Wang He fue la estimación de pose a nivel de categoría: ante una taza que nunca ha visto, ¿puede el modelo saber que la boca de la taza apunta hacia arriba y que el asa está a la derecha, sin tener que escanear y modelar cada taza ni definir un sistema de coordenadas para cada una?

Los datos reales prácticamente no existían. Las imágenes de internet no tienen anotaciones de pose en seis dimensiones, y un estudiante de doctorado no puede fotografiar ni anotar una cantidad ilimitada de tazas. Por eso recurrió a datos de graficación por computadora y de realidad mixta: fue a IKEA a fotografiar fondos reales de mesas con su profundidad, renderizó tazas digitales sobre esos fondos reales, obtuvo automáticamente las etiquetas de pose y generó cientos de miles de imágenes de entrenamiento. Este método coincide en gran medida con las ideas actuales de datos sintéticos y Sim2Real: el fondo real aporta anclaje en el mundo real, el primer plano virtual aporta etiquetas controlables y, al final, la evaluación se hace en el mundo real.

\begin{importantbox}{Una versión temprana de los datos sintéticos}
La estimación de pose a nivel de categoría muestra que los datos sintéticos no son una palabra de moda surgida apenas hoy. Siempre que la anotación de datos reales no puede llevarse a gran escala, los investigadores recurren a la graficación por computadora, la simulación o la realidad mixta para fabricar muestras de entrenamiento controlables. Los datos sintéticos actuales para robots solo extienden esta idea a más objetos, acciones, escenas e interacciones físicas.
\end{importantbox}

\subsection{Por qué insistir en los robots domésticos}

Esta subsección lleva la trayectoria personal hacia las decisiones estratégicas. Alrededor de 2020, Li Kaifu (李开复) le sugirió a Wang He aplicar sus capacidades en visión tridimensional y estimación de pose a la conducción autónoma, el LiDAR y la estimación de pose de vehículos; Wang He, en cambio, insistió en trabajar en home robot, porque consideraba que el entorno doméstico ofrece el espacio de interacción más amplio y es el más cercano a un «agente inteligente general en el mundo físico». En ese momento era una apuesta muy a contracorriente: se pensaba que los robots domésticos todavía estaban muy lejos, y en China casi no había aliados.

Tras volver a China, Wang He impulsó la línea de inteligencia corpórea en la Universidad de Pekín (北大) y en el BAAI (智源), escribió artículos para presentar la embodied AI y, junto con unos pocos investigadores como Lu Cewu (卢策武), introdujo este concepto en la discusión académica e industrial del país. El punto didáctico aquí no es la leyenda personal, sino la estructura de una elección temprana de dirección: si un investigador cree en la next wave, tiene que soportar un periodo sin consenso, sin aliados y sin interés del capital.

\subsection{Resumen de la sección}

La trayectoria académica de Wang He dejó al descubierto de antemano dos problemas centrales de la inteligencia corpórea: primero, la inteligencia debe comprender cómo las acciones cambian el mundo; segundo, la falta de datos reales debe suplirse con datos sintéticos o de realidad mixta. El objetivo de los robots domésticos, por su parte, lleva estos problemas de investigación hacia aplicaciones más amplias en el mundo físico.
\section{Ascenso en espiral del software y el hardware, y la brecha de datos}

La primera mitad trató la historia académica; esta sección entra en la ruta de industrialización. Wang He (王鹤) sostiene que el software y el hardware de la inteligencia corporizada forman un problema de ascenso en espiral: un hardware demasiado agresivo e inmaduro frena el entrenamiento de la inteligencia y su implantación en escenarios reales; a su vez, cuando mejoran las capacidades del software, estas exigen un hardware más adecuado para la recolección de datos, la ejecución y la entrega. No se puede concebir el robot como «primero construir un cuerpo general y luego esperar a que el cerebro lo aprenda todo», ni tampoco hacer solo software sin atender las restricciones del cuerpo físico.

La descripción de la entrevista contiene un juicio central: el valor de producción global de todos los brazos robóticos industriales el año pasado fue de apenas unos 100,000 millones de RMB, comparable al de una sola automotriz como Lixiang (理想). El significado de esta cifra es que el mercado comercial real de la industria robótica todavía no es grande, y no se puede valuar solo con narrativas grandilocuentes. La ruta de Yinhe Tongyong (银河通用) consiste en desarrollar una inteligencia relativamente especializada sobre el hardware existente y avanzar de forma gradual hacia una inteligencia cada vez más general. Aquí «especializada» no significa renunciar a lo general, sino encontrar primero escenarios capaces de generar productividad.

\begin{figure}[H]
\centering
\includegraphics[width=0.96\textwidth]{software-hardware-spiral.png}
\caption{Ascenso en espiral del software y el hardware: un hardware demasiado agresivo frena la inteligencia, y una inteligencia madura, a su vez, impulsa el hardware. Diagrama conceptual propio, elaborado a partir de la conversación de 01:25:08--01:44:34 y la descripción del video.}
\end{figure}

\begin{knowledgebox}{Lectura de la figura: la espiral no es un pipeline lineal}
El hardware, el cuerpo físico, los datos, el modelo y el escenario no son módulos que se completan uno tras otro, sino que se condicionan mutuamente. El hardware determina qué datos se pueden recolectar y qué acciones se pueden ejecutar; el modelo determina qué sensores y qué espacio de acciones se necesitan; el escenario determina los requisitos de costo y confiabilidad. Al leer la figura hay que fijarse en el ciclo, en lugar de tratar el hardware como una compra única.
\end{knowledgebox}

\subsection{Por qué los VLM son más débiles que los LLM}

Esta sección explica una brecha de datos fundamental. Los datos de texto de los LLM se acercan a una cobertura a gran escala de «lo que los seres humanos han dicho en internet»; los datos visuales de los VLM cubren solo una pequeña parte de «lo que el ojo humano observa del mundo»; los VLA, además, deben incorporar action data, y la recolección sistemática de action data comenzó apenas en los últimos dos años. Dicho de otro modo, los modelos de visión y lenguaje no son inferiores por una arquitectura intrínsecamente más primitiva, sino porque su cobertura de datos, su estructura de anotación y su retroalimentación de acciones son mucho más débiles que las del texto.

Esto explica por qué los robots no pueden simplemente reutilizar la scaling law de los modelos de lenguaje. Un modelo de lenguaje se entrena con secuencias de tokens, e internet ya proporciona enormes cantidades de texto; un robot se entrena con el acoplamiento entre visión, lenguaje, acción, estado, fallo, recuperación y escenario, y esos datos son por naturaleza más escasos, más caros y más sesgados. Los datos sintéticos, la teleoperación, la simulación y los datos que regresan de las unidades entregadas son importantes precisamente porque intentan cubrir el déficit de action data.

\begin{knowledgebox}{Términos clave: diferencias de datos entre LLM, VLM y VLA}
\begin{tabularx}{\textwidth}{p{0.18\textwidth}p{0.34\textwidth}X}
\toprule
Tipo de modelo & Cobertura de datos & Limitación principal \\
\midrule
LLM & El texto de internet se acerca a cubrir una gran parte de la expresión humana y del intercambio de conocimiento & El texto no contiene directamente el costo de la ejecución física. \\
VLM & Las imágenes y los videos cubren solo una pequeña parte de la experiencia visual humana & Las observaciones carecen de retroalimentación tridimensional, temporal y causal completa. \\
VLA & Además requiere trayectorias de acción, estados de las articulaciones o del efector final y resultados de ejecución & La recolección sistemática de action data apenas comienza y su costo real es altísimo. \\
Política robótica & Requiere acoplar visión, lenguaje, acción, recuperación ante fallos y objetivos del escenario & Los datos deben estar anclados al hardware y a tareas reales. \\
\bottomrule
\end{tabularx}
\end{knowledgebox}

El ponente enfatiza que no hay que fijar una meta demasiado alta de golpe: construir un VLA completamente general en uno o dos años no es realista, dadas las condiciones actuales de datos e industria. Una ruta más sólida es desarrollar la inteligencia en torno a una application replicable a gran escala, de modo que el modelo cubra todo lo que es necesario hacer dentro de esa aplicación. Por ejemplo, en escenarios de supermercados y farmacias, el desplazamiento, la sujeción, la colocación, el reabastecimiento y el acomodo forman un conjunto de habilidades limitado, pero los objetos, las marcas, las tiendas y las formas de exhibición deben generalizarse por completo.

\begin{figure}[H]
\centering
\includegraphics[width=0.94\textwidth]{vlm-vs-llm-data.png}
\caption{Por qué los VLM son más débiles que los LLM: la cobertura visual y los datos de acción son mucho menores que la cobertura textual. Diagrama conceptual propio, elaborado a partir de la conversación de 01:25:08--01:44:34 y la descripción del video.}
\end{figure}

\begin{knowledgebox}{Lectura de la figura: la brecha está en la cobertura de datos, no solo en el nombre del modelo}
A la izquierda, los LLM disfrutan de una cobertura textual a gran escala; a la derecha, los VLM y VLA necesitan visión, información tridimensional, acciones y retroalimentación. En los VLA, la capa de Action Data es la más escasa: sin trayectorias de acción ni resultados, al modelo le resulta muy difícil pasar de «entender lo que ve» a «lograr hacerlo».
\end{knowledgebox}

\begin{warningbox}{Una propuesta de hardware agresiva frena la inteligencia}
Si el hardware es inmaduro, su mantenimiento es costoso y su espacio de acciones es inestable, el entrenamiento del modelo tendrá que adaptarse una y otra vez a un cuerpo físico defectuoso, y la entrega en los escenarios quedará absorbida por los problemas de confiabilidad. La inteligencia corporizada no es mejor cuanto más general sea; se trata de encontrar, dentro de lo que el hardware actual puede sostener, una inteligencia capaz de generar valor real.
\end{warningbox}

\subsection{Resumen de la sección}

La inteligencia corporizada no es un problema puramente de software ni puramente de hardware. El software, el hardware, los datos y los escenarios deben ascender en espiral; la brecha de los VLM y VLA frente a los LLM radica, en esencia, en la cobertura visual insuficiente y en la escasez de action data.
\section{Dos atolladeros: flotar indefinidamente y que las cuentas no cuadren}

La sección anterior trató el software, el hardware y los datos; esta sección aborda los riesgos de las empresas emergentes. Wang He (王鹤) sostiene que, si una empresa de inteligencia corporizada cae en alguno de dos atolladeros, su techo será muy limitado. El primero es «flotar indefinidamente»: la empresa permanece mucho tiempo en el terreno de los conceptos, las demostraciones y el relato para atraer financiamiento, sin entrar en escenarios reales ni someterse a la prueba de los clientes, los costos y la confiabilidad. El segundo es que «las cuentas no cuadren»: aun cuando entra en un escenario, el costo marginal no baja, cada entrega pesa más que la anterior y, al final, la productividad no se materializa.

En esencia, ambos atolladeros son un desajuste en la retroalimentación. Flotar indefinidamente significa que el reward model de la empresa es la attention y la valuation; que las cuentas no cuadren significa que la ruta tecnológica de la empresa no ha logrado reducir el costo unitario de entrega. La inteligencia corporizada es una industria física: no basta con mirar el benchmark del modelo ni los videos demo. Tiene que responder lo siguiente: una vez que un robot entra en un escenario, ¿puede completar un tipo de tarea de manera más estable, más barata o más escalable que una persona?

\begin{figure}[H]
\centering
\includegraphics[width=0.94\textwidth]{two-traps.png}
\caption{Dos atolladeros: tanto flotar indefinidamente como que las cuentas no cuadren limitan el techo de la empresa. Diagrama conceptual de elaboración propia, a partir de la conversación entre 01:44:34 y 01:55:17 y de la descripción del video.}
\end{figure}

\begin{knowledgebox}{Lectura de la figura: uno es un riesgo de relato y el otro, un riesgo económico}
El problema de flotar indefinidamente es desconectarse de los clientes reales; el problema de que las cuentas no cuadren es que cuanto más se entrega, más se pierde. Al leer la figura hay que relacionar ambos lados: hacer solo demos evita exponer los costos a corto plazo, pero también hace perder los datos reales; y si, una vez dentro del escenario, el costo marginal no baja, el volante de datos tampoco puede sostener el ciclo comercial cerrado.
\end{knowledgebox}

\subsection{Generalización dentro del escenario de aplicación}

Esta sección explica por qué Galbot, Yinhe Tongyong (银河通用), eligió la «generalización dentro del escenario de aplicación». Los robots de propósito general resultan atractivos, pero buscar demasiado pronto la generalización a todos los escenarios hace que las tareas, el hardware, los datos y los costos se disparen al mismo tiempo. Una ruta más realista es generalizar a fondo dentro de un escenario de aplicación concreto: por ejemplo, en el escenario de anaqueles, los distintos productos, su acomodo, las oclusiones, los faltantes, el reabastecimiento y las situaciones anómalas. No se trata de automatizar un solo punto, sino de construir capacidades escalables dentro de un mundo acotado.

Este enfoque se parece también a la conversión en producto de los modelos grandes: primero se logra que funcionen de punta a punta el dominio de tareas, la retroalimentación, los costos y el ciclo cerrado de datos, y luego se amplían gradualmente los límites de la capacidad. Para los robots, la generalización dentro del escenario tiene una ventaja adicional: una vez que hay unidades entregadas, los datos reales regresan y el sistema puede aprender de las fallas reales. Sin volumen de entregas, el volante de datos es solo una ilusión.

\begin{importantbox}{Generalización dentro del escenario}
En la etapa temprana de la inteligencia corporizada, la ruta eficaz quizá no sea «hacer primero un robot de propósito general», sino lograr una generalización suficiente de objetos, tareas, anomalías y costos dentro de un escenario real por el que se paga. Cuanto más claros sean los límites del escenario, más fácil será validar el ciclo cerrado de datos y el ciclo comercial cerrado.
\end{importantbox}

\subsection{Resumen de la sección}

Las empresas de inteligencia corporizada deben evitar dos atolladeros: no quedarse flotando indefinidamente en el relato y no entrar en un escenario donde las cuentas no cuadren. La ruta correcta es encontrar un escenario por el que se pague, generalizar dentro de él y lograr que las unidades entregadas devuelvan datos.
\section{La productividad es el producto: volante de datos y entrega real}

La sección anterior explicó que la empresa debe evitar dos atolladeros, el del relato y el de la rentabilidad; esta sección responde una pregunta más: si no se apoya en una historia, ¿qué vende en realidad la inteligencia corporizada? Los modelos grandes pueden decir «la inteligencia es el producto», porque el usuario compra directamente capacidades de conversación, redacción, código, búsqueda y Agent; la inteligencia corporizada, en cambio, tiene que decir «la productividad es el producto». Un robot no solo exhibe inteligencia: tiene que crear en el mundo físico una sustitución de trabajo, una mejora de eficiencia o una capacidad productiva nueva que puedan medirse. Lo que el cliente compra al final no es que «se parezca a una persona», sino la producción estable de tareas como acomodar anaqueles, cargar y trasladar, seleccionar artículos, inspeccionar, limpiar y repartir.

Esto también explica por qué la recolección de datos reales es tan cara. Contratar personas para teleoperar robots y recolectar datos reales tiene un costo alto, avanza despacio y cubre pocos escenarios; la descripción del episodio menciona que los datos reales representan solo el 1\% de los datos de entrenamiento y que el pipeline de datos sintéticos lleva el peso principal. Esa proporción no puede trasladarse mecánicamente a todas las empresas, pero muestra que, sin datos sintéticos, simulación y generalización dentro del escenario, el costo de los datos reales difícilmente sostendría la escala temprana de la inteligencia corporizada.

\begin{figure}[H]
\centering
\includegraphics[width=0.96\textwidth]{productivity-as-product.png}
\caption{La productividad es el producto: los modelos grandes venden inteligencia; la inteligencia corporizada tiene que vender productividad medible. Diagrama conceptual propio, elaborado a partir de la conversación de 01:55:17--02:13:51 y la descripción del video.}
\end{figure}

\begin{knowledgebox}{Lectura de la figura: de la tarea al ingreso hay que pasar por el costo}
En la figura, el paso de Task a Robot es solo el inicio; lo decisivo son Cost, Output y Revenue. Para que un robot genere valor de producto, tiene que bajar el costo por tarea, entregar una producción estable y verificable, y lograr que el cliente quiera volver a comprar. Una inteligencia que no reduce costos es solo una demostración cara.
\end{knowledgebox}

\subsection{Las cuentas de la teleoperación}

Esta subsección pone en números la afirmación «los datos reales son caros». La estimación que da Wang He (王鹤) es la siguiente: fabricar un robot humanoide full-size cuesta por lo menos del orden de cien mil yuanes; si se compran diez mil unidades para recolectar datos, la inversión en hardware es del orden de mil millones de yuanes. Más cara aún es la operación: si cada robot trabaja en dos turnos, incluso con dos personas teleoperándolo por turno, y se suman la anotación, el control de calidad y el mantenimiento, sostener durante un mes un sistema de recolección de diez mil unidades cuesta del orden de cientos de millones e incluso mil millones de yuanes. Estas cuentas explican por qué «depender por completo de datos reales de teleoperación» difícilmente puede ser la ruta general en la etapa temprana.

\begin{knowledgebox}{Las cuentas: por qué no basta con teleoperar para recolectar datos}
\begin{tabularx}{\textwidth}{p{0.22\textwidth}p{0.28\textwidth}X}
\toprule
Rubro de costo & Orden de magnitud aproximado & Lección \\
\midrule
Hardware de los robots & Diez mil unidades $\times$ del orden de cien mil yuanes & Antes de recolectar ya se requiere una inversión en activos del orden de mil millones de yuanes. \\
Personal de teleoperación & Varios turnos y varias personas por unidad & El personal no es un costo único, sino un burn continuo. \\
Anotación y control de calidad & Hay que revisar trayectorias, estados, fallas y resultados de la tarea & Los datos de acción no pueden etiquetarse a grandes trazos como una imagen común. \\
Escenario y mantenimiento & Instalaciones, reparaciones, reabastecimiento, manejo de excepciones & El mundo físico produce sin cesar problemas que no son del modelo. \\
\bottomrule
\end{tabularx}
\end{knowledgebox}

Esta es también una diferencia importante entre la conducción autónoma y la inteligencia corporizada. Un auto puede venderse al usuario, y el usuario genera datos de paso al manejarlo todos los días, de modo que el uso del producto absorbe el costo de los datos; si un robot todavía no puede crear valor por sí mismo, no hay manera de disfrazar la recolección de datos como comercialización. Vender al cliente un «robot sin funciones» para que él mismo recolecte y entrene puede generar ingresos a corto plazo, pero a largo plazo traslada al cliente el riesgo de investigación y desarrollo.

\subsection{La recipe de datos y el retorno desde las unidades entregadas}

Esta subsección convierte el problema de los datos en un problema de volante. La recipe de datos de la inteligencia corporizada puede entenderse así: los datos de internet y visuales aportan sentido común y representaciones; los datos sintéticos aportan la cola larga y un entrenamiento controlable; una cantidad pequeña de datos reales aporta grounding, y las fallas reales después de la entrega aportan un retorno continuo. En la etapa temprana, sin volumen de unidades entregadas, la proporción de datos reales puede ser muy baja; en cuanto se entra a escenarios reales, el retorno de datos conecta el producto con el pipeline de entrenamiento.

\begin{figure}[H]
\centering
\includegraphics[width=0.92\textwidth]{data-recipe-robotics.png}
\caption{Recipe de datos para sistemas corporizados: 1\% de datos reales, pipeline sintético, retorno desde las unidades entregadas y volante de datos. Diagrama conceptual propio, elaborado a partir de la conversación de 01:55:17--02:13:51 y la descripción del video.}
\end{figure}

\begin{knowledgebox}{Lectura de la figura: cuanto más arriba, más cerca de la realidad y más escaso}
La capa inferior de datos generales construye las capacidades básicas; la capa intermedia de datos sintéticos amplía la distribución de tareas; una cantidad pequeña de datos reales calibra, y el retorno desde las unidades entregadas aporta los casos de falla y de frontera más importantes. Al leer la figura hay que evitar un malentendido: que los datos reales sean una proporción baja no significa que sean poco importantes; precisamente porque son escasos, deben usarse para calibrar y validar todo el sistema.
\end{knowledgebox}

\begin{warningbox}{Un robot sin funciones no puede venderse como volante de datos}
La descripción de la entrevista menciona un fenómeno tricky: vender a otros un robot sin funciones, dejar que el cliente recolecte y entrene por su cuenta, y prometerle que en el futuro se convertirá en un empleado. Este modelo es muy riesgoso, porque traslada al cliente el riesgo de investigación y desarrollo y además agota la credibilidad de la industria.
\end{warningbox}

\begin{importantbox}{La definición de productividad}
Wang He define la productividad de forma muy sencilla: el trabajo realizado por unidad de tiempo debe equipararse al de una persona o, por lo menos, mejorar la eficiencia dentro de un límite claro. Si el robot es mucho más lento que una persona, no aguanta mucho tiempo, falla con frecuencia o necesita mucho respaldo humano, no es una productividad de calidad; incluso puede ser una productividad atrasada.
\end{importantbox}

\subsection{Resumen de la sección}

El producto de la inteligencia corporizada no es «parecer inteligente», sino productividad medible. El volante de datos también tiene que apoyarse en entregas reales: los datos sintéticos reducen el costo inicial, el retorno desde las unidades entregadas aporta calibración real y el valor para el cliente comprueba si la ruta funciona.
\section{El desorden provocado por la avalancha de capital}

La sección anterior trató la productividad; esta trata el desorden en la industria. Después de ChatGPT, la inteligencia corporizada pasó de ser un concepto minoritario a convertirse en la favorita del capital, y el financiamiento, las valuaciones y los demo se dispararon. El núcleo de la crítica de Wang He (王鹤) es este: el mayor desorden de la industria consiste en no poder distinguir quién crea productividad y quién solo cuenta una historia. Más grave aún, algunos equipos no aclaran en sus presentaciones que usan teleoperación, e incluso hacen pasar la teleoperación por capacidad autónoma, lo que destruye directamente la confianza del público y de los clientes en la industria.

Este tipo de problema no es solo ético; también se convierte en un problema técnico. Las presentaciones engañosas desvían el reward model de la industria: los equipos optimizan los demo para conseguir financiamiento y difusión, en lugar de optimizar el sistema para la tasa de éxito en escenarios reales, la reducción de costos y el retorno de datos de las unidades entregadas. A largo plazo, el capital castigará a toda la industria, y los equipos que de verdad generan productividad también se verán arrastrados por el ruido.

\begin{figure}[H]
\centering
\includegraphics[width=0.92\textwidth]{capital-chaos-reward.png}
\caption{El desorden después de la avalancha de capital: quién crea productividad y quién solo cuenta una historia. Figura conceptual propia, elaborada a partir de la conversación de 02:13:51--02:25:25 y de la descripción del video.}
\end{figure}

\begin{knowledgebox}{Lectura de la figura: un reward equivocado produce conductas equivocadas}
En el centro de la figura está el reward. Si el reward proviene del monto de financiamiento, la valuación y el tráfico, los equipos optimizarán naturalmente la historia y la demostración; solo si el reward proviene de presentaciones veraces, de la recompra de los clientes y de aplicaciones de diez mil unidades, los equipos optimizarán la productividad. Al leer la figura, conviene considerar la «presentación veraz» como infraestructura de la industria: sin presentaciones confiables, el mercado no puede juzgar el progreso.
\end{knowledgebox}

\subsection{El reloj de validación: diez mil unidades en cinco años}

Wang He plantea un juicio contundente: en un plazo de cinco años, el campo debe contar con aplicaciones de más de diez mil unidades; si no lo logra, el campo de la inteligencia corporizada quedará refutado, o al menos lo quedará su narrativa actual. Esta afirmación no es necesariamente una predicción precisa, pero ofrece un criterio de aceptación para la industria. Las diez mil unidades no buscan una cifra vistosa, sino demostrar que el hardware, el software, la entrega, el mantenimiento, el retorno de datos y el valor para el cliente ya superaron la etapa de laboratorio.

\begin{figure}[H]
\centering
\includegraphics[width=0.96\textwidth]{robotics-validation-clock.png}
\caption{Reloj de validación de cinco años: se necesitan aplicaciones de más de diez mil unidades en cinco años; de lo contrario, la narrativa del campo quedará refutada. Figura conceptual propia, elaborada a partir de la conversación de 02:13:51--02:25:25 y de la descripción del video.}
\end{figure}

\begin{importantbox}{Principio de presentación veraz}
Se puede usar teleoperación, pero hay que decirlo; se pueden mostrar productos inacabados, pero hay que aclarar sus límites; se puede buscar financiamiento, pero no se puede sustituir el progreso real con demo irreproducibles. Para la inteligencia corporizada, la credibilidad de la industria es en sí misma un bien público.
\end{importantbox}

\begin{knowledgebox}{Matriz de validación: qué cuenta como progreso real}
\begin{tabularx}{\textwidth}{p{0.20\textwidth}p{0.34\textwidth}X}
\toprule
Elemento de validación & Qué hay que observar & Con qué no puede sustituirse \\
\midrule
Presentación pública & En vivo, sin teleoperación oculta, tareas reproducibles & Videos editados y fragmentos de un único intento exitoso. \\
Operación en escenarios & Cuántos pedidos o tareas reales se procesan al día y durante cuánto tiempo & Alianzas estratégicas y cartas de intención. \\
Escala de entregas & Si se alcanzan despliegues reales del orden de miles o decenas de miles de unidades & El número de prototipos de ingeniería. \\
Productividad & Producción por unidad de tiempo, tasa de fallas, proporción de respaldo humano & Apariencia y movimientos «parecidos a los humanos». \\
Volante de datos & Si las fallas del despliegue regresan al entrenamiento y a la mejora del producto & Centros de recolección estáticos y demostraciones únicas. \\
\bottomrule
\end{tabularx}
\end{knowledgebox}

Esta matriz también explica por qué Wang He insiste en no dañar la reputación de la industria. China enfrenta el envejecimiento de su población y un déficit de mano de obra; la productividad robótica no es solo una oportunidad de emprendimiento, sino que podría influir en la oferta a largo plazo de la manufactura y los servicios. Si la industria se ve arrastrada a una «era de hielo» por presentaciones falsas y promesas poco realistas, el problema social de quienes realmente necesitan robots para complementar la mano de obra también perderá una posible solución.

\subsection{El significado del episodio con Jensen Huang}

Lo anterior trató la credibilidad de la industria; esta sección usa el episodio final de la entrevista para volver al ecosistema tecnológico. Durante la visita de Jensen Huang (黄仁勋) a China, Wang He compartió mesa con él y se sentó a su lado; observó que Jensen Huang tolera el picante, que comió bastante carne de cerdo hervida en aceite picante y que respondió con entusiasmo a una función de cambio de máscaras. El fragmento parece anecdótico, pero en estas notas puede servir como indicio indirecto: NVIDIA, la IA física, la inteligencia corporizada y los emprendedores chinos de robótica ya se encuentran en el mismo campo industrial. El cómputo, la simulación, el cuerpo del robot y los escenarios de aplicación ya no son líneas paralelas.

La explicación que da Wang He de esa comida compartida también se relaciona con el eje técnico: NVIDIA concede mucha importancia a los datos sintéticos, porque estos requieren renderizado, simulación y entrenamiento en GPU; si esta ruta funciona, las empresas de GPU podrán sostener la mitad de la inteligencia corporizada. Galbot (银河通用) también mostró a NVIDIA su trabajo de entrenamiento con datos sintéticos y transferencia al mundo real. Esto muestra que «la productividad es el producto» y «los datos sintéticos son la clave de la recipe de datos» no son dos juicios separados, sino los dos extremos de un mismo ciclo cerrado: la producción de datos sostiene la capacidad, la capacidad entra en los escenarios y genera productividad, y los escenarios devuelven datos reales.

\subsection{Resumen de la sección}

El capital amplifica las oportunidades, pero también el desorden. Si la inteligencia corporizada quiere construir credibilidad a largo plazo, debe desplazar el reward del financiamiento y los demo hacia las presentaciones veraces, la productividad, la escala de entregas y el valor para el cliente.
\section{Síntesis y ampliación}

Esta sección condensa toda la entrevista en cinco conclusiones. Primera: en su origen, la inteligencia corporizada no fue una extensión natural de la industria robótica tradicional, sino una migración académica de la visión por computadora, que pasó del reconocimiento pasivo a la acción activa. Segunda: el lenguaje no es la esencia de la inteligencia, sino un salto en la inteligencia humana; la inteligencia robótica debe volver a la reacción ante el entorno y a la retroalimentación física. Tercera: el cuello de botella de los VLM/VLA se debe en gran medida a una cobertura de datos insuficiente, en particular a la escasez de action data. Cuarta: software, hardware, datos y escenarios de aplicación deben avanzar juntos en espiral; no se puede apostar por separado a un cuerpo robótico de propósito general ni a un cerebro de propósito general. Quinta: la industria, a fin de cuentas, tiene que demostrar su valor con productividad, unidades entregadas y demostraciones reales.

\begin{importantbox}{EP106 dentro de la serie de Zhang Xiaojun (张小珺) sobre robótica}
EP106 aporta la historia académica y la ética del sector; EP109, la infraestructura de simulación y datos sintéticos; EP121, la perspectiva de los modelos robóticos de DeepMind, y EP132, casos de robots completos y de cadena de suministro. En conjunto, los cuatro episodios forman una ruta de aprendizaje sobre robótica: «origen del concepto--producción de datos--ruta de modelos--implementación industrial».
\end{importantbox}

\subsection{Takeaways clave}

\begin{enumerate}
\item Lo esencial de la inteligencia corporizada no es la forma del robot, sino el ciclo cerrado de percepción, acción y retroalimentación del entorno.
\item El lenguaje potencia la inteligencia, pero la inteligencia del mundo físico debe anclarse en el cuerpo y en los action data.
\item Los datos reales son muy caros; los datos sintéticos y los datos que regresan de las unidades entregadas son componentes clave de la recipe de datos para robots.
\item Las empresas en etapa temprana deben evitar quedarse a la deriva por mucho tiempo y operar con números que no cuadran; primero deben lograr la generalización dentro de un escenario de aplicación.
\item La credibilidad del sector proviene de demostraciones reales y de la validación en productividad; la teleoperación, el financiamiento y las demo no sustituyen el avance real.
\end{enumerate}

\subsection{Lecturas adicionales}

\begin{itemize}
\item Para entender la infraestructura de simulación y datos sintéticos, puede contrastarse con la entrevista a Xie Chen (谢晨) en EP109.
\item Para entender los modelos base para robots, la transferencia entre distintos cuerpos robóticos y los modelos del mundo, puede contrastarse con la entrevista a Tan Jie (谭捷), de DeepMind, en EP121.
\item Para entender los robots completos, la cadena de suministro y la Data Recipe, puede contrastarse con la entrevista a Gao Jiyang (高继扬), de Xinghaitu (星海图), en EP132.
\end{itemize}

\end{document}
